\section{One-Way to Hiding, Revisited}\label{sec:improving-o2h}
  Now that we have compressed random oracles in our toolbox, it is natural to
    ask if we can somehow use the to derive a better One-Way to Hiding bound.
    After all, the lossiness of the O2H lemmas stand in stark contrast to the
    classical case, but without appearing to reflect any real-world attacks.

  \hhnote{
    Think about whether there is an ``obvious'' application of Grover that I'm
    missing, that the square root could be attributed to. \\
    A: Seems like there might be, given the example proof in
    \cite{C:AmbHamUnr19}!
  }

  While a full impossibility theorem for tighter bounds is missing, the
    following two lemmas provide some insight into why we should expect
    1.\ the square root, and 2.\ some loss factor $c(q_{\RO})$
    to be inevitable.

  The first lemma takes a ``gray-box'' view of the adversary, and shows that
    there is a separation between observing a value and distinguishing two
    oracles, which seems to be unique to the quantum world.

  \begin{lemma}
    Let $\adv$ get an input $z$ that depends arbitrarily on $\realRO(x)$ for $x
      \in \careset$, and let
      $\idealRO$ be an independent random oracle. Given oracle access to
      $\bitRO$, $\adv$ is tasked with outputting $\chalbit$. Let $\adv$
      also output a second bit encoding the event $\obsev$, set to
      $1$ if and only if $\adv$ ``observed'' a value $\bitRO$ for $x \in
      \careset$. Then for all such $\adv$ we have, 
      \[ 
        \prob{\obsev} \leq 
          \advantage{\hide}{\bitRO, z}[\distinguisher] 
          \leq \prob{\obsev} 
          + 2 \sqrt{ 
            2 \cdot \prob{\obsev} \left( 1 - \prob{\obsev} \right)
          }\, .
      \]
      Equality on the right may be achieved by applying a controlled $\Xgate$
      gate to the ancilla $\qubit_\obsev$, and equality on the left is achieved by
      applying a controlled $i \Xgate$ gate instead.
  \end{lemma}

  \noindent
  The right-hand-side in the above lemma may of course be further bounded to give the
  more familiar-looking
  \[ 
    %\prob{\obsev} + 2 \sqrt{ 
    %  2 \cdot \prob{\obsev} \left( 1 - \prob{\obsev} \right)
    %} 
    \advantage{\hide}{\bitRO, z}[\distinguisher] 
      \leq 2 \sqrt{2 \cdot \prob{\obsev}}\, .
  \]

  \begin{proof}
    \hhnote{to-do: adapt the Hands-Off O2H proof}
    \qed
  \end{proof}

  \noindent
  Now, if an extractor $\extractor$ were lazy-sampling a classical random
    oracle, then the event $\obsev$ would correspond exactly to $\extractor$
    being able to extract. Even if we accept
    that a square-root loss seems inevitable, can we at least get rid of the
    factors $q_{\RO}$ by simulating the random oracle as a CRO, letting $\adv$
    run until the end, and measuring the database to see what we get?

  Unfortunately, the following lemma answers also this in the negative.

  \begin{lemma}\label{lemma:uncomputing-adv}
    Let $z$ be some value that depends arbitrarily on $\realRO(x)$ for $x
    \in \careset$. Then there is a distinguisher $\distinguisher$ and an value
    $z$, such that when $\distinguisher$ is given $z$,
    \[
      \advantage{\hide}{\idealRO, z}[\distinguisher^{\realRO}] = 1\, ,
      \qquad \advantage{\hide}{\idealRO, z}[\distinguisher^{\idealRO}] > 1 - \frac{1}{2^m}\, ,
    \]
    and yet for any extractor $\extractor$ that simulates
      $\idealRO$ as a compressed random oracle, without
      ever aborting $\adv$, and observes the oracle database at the end of the
      protocol,
    \[
      \advantage{\oneway}{}[\extractor] < \frac{1}{2^m}\, ,
    \]
    where $2^m$ is the size of the random oracles' output domains.
  \end{lemma}
  
  \begin{proof}
    Let $\adv$ be as given in \autoref{fig:uncomputing-adv} \ldots
    \hhnote{
      To-do: Add simple proof using gentle operator lemma.
    }

    \begin{figure}[t]
      \centering
      \begin{pchstack}[center]
        \procedure{
          Adversary $\adv^{\idealRO}(\target, \targetout \coloneqq \realRO(\target))$
        }{
          \bitreg \concat \inputreg \concat \outputreg \gets \reginit{(1+n+m)}{\ket{0}} \\
          \ket{\inputreg} \gets \ket{\target} \\
          (\inputreg, \outputreg) \gets \idealRO(\inputreg, \outputreg) \\
          \pcforeach \ket{y} \pcin\ \outputreg \pcdo \\
            \t \pcif y = \targetout \pcthen \applyU{\Xgate}{\bitreg} \\
          (\inputreg, \outputreg) \gets \idealRO(\inputreg, \outputreg) \\
          \pcreturn \register[B]
        }

        \pchspace

        \procedure{Oracle $\idealRO(\inputreg, \outputreg : \database)$}{
          \applyU{\decomp}{\inputreg \concat \database} \\
          \applyU{\Hgate^{\otimes m}}{\outputreg} \\
          \applyU{\entangle}{\inputreg \concat \outputreg \concat \database} \\
          \applyU{\Hgate^{\otimes m}}{\outputreg} \\
          \applyU{\decomp}{\inputreg \concat \database} \\
          \pcreturn (\inputreg, \outputreg)
        }
      \end{pchstack}
      \caption{
        \label{fig:uncomputing-adv}
        \emph{Left:} An adversary $\adv$ that after an RO query, entangles to an output
        qubit, and then
        uncomputes the query, in a scenario where it is given an
        input corresponding to a query to the ``real'' oracle $\realRO$, while
        given oracle access to the independent, simulated oracle $\idealRO$.
        $\adv$ is then tasked with noticing the difference, which it can
        trivially do by calling $\idealRO$ on $\target$. \\
        %Here, $\inputreg$, $\outputreg$, and $\bitreg$ are $n$-, $m$-, and
        %$1$-length qubit registers, respectively, initialized to the all-zero
        %states. \\
        \emph{Right:} The simulated oracle $\idealRO$, implemented as a compressed
        oracle (taken from \autoref{def:compressed-ros}). Here, $\database$ is
        the database, initialized to the list of $q$ identical tuples
        $(\emptystring, 0^m)$, (i.e.\ at the outset the method $\database \gets
        \reginit{q}{\ket{\emptystring, 0^m}}$ is run). 
        %Since there are only $2$ queries in this example, we may set $q=2$.
      }
    \end{figure}

    \qed
  \end{proof}

  \noindent
  The main difficulty is that $\adv$ can uncompute queries after entangling
    with them simply by repeating the query with the same registers. If the
    entanglement doesn't perturbe the state too
    much, then $\database$ will also be ``almost unperturbed'', meaning it will
    have a very low amplitude of storing the input value.
    In order to avoid this uncomputing, it appears that we have no choice but
    to let the extractor abort the distinguisher, either (as in the original O2H) at a random point,
    or (as in the Semi-Classical O2H) by checking each query for a bad event.
    Regardless of the approach taken, this is guaranteed to lead to a loss
    factor that depends on $q_{\RO}$.

  Another way to prove the \autoref{lemma:uncomputing-adv} is to do a direct calculation of the final
    state of the combined system, which amounts to an instructive exercise
    in reasoning about compressed oracles.

  \begin{exercise}
    Let $\adv$ perform the operations given in \autoref{fig:uncomputing-adv}
    (left), and let $\extractor$ simulate $\idealRO$ as a compressed random oracle
    (\autoref{fig:uncomputing-adv}, right). Denoting the
    the registers holding the first
    tuple of the oracle database as $\database_0$, show that the final 
    state of the combined system before measuring is
    \begin{align*}
      \ket{\bitreg}\ket{\inputreg}\ket{\outputreg}\ket{\database_0}
      &= \left(
          \left(
            1 - \frac{1}{2^m}
          \right) \ket{0}_{\bitreg} 
          + \frac{1}{2^{m}}
          \ket{1}_{\bitreg} 
        \right)
        \ket{\target}_{\inputreg} 
        \ket{0^m}_{\outputreg}
        \ket{\emptystring, 0^m}_{\database_0} \\
      &\qquad +
        \frac{1}{2^{m}}
        \sum_{u \neq 0^m}
        (-1)^{\targetout \cdot u}
        \left(
          \ket{1}_{\bitreg} 
          - \ket{0}_{\bitreg} 
        \right)
        \ket{\target}_{\inputreg} 
        \ket{0^m}_{\outputreg}
        \ket{\target, u}_{\database_0}\, ,
    \end{align*}
    and show that observing this state in the computational basis yields the
      bounds stated in \autoref{lemma:uncomputing-adv}.
      You will need the following relations: The Hadamard transformation,
    \begin{align*}
      \Hgate^{\otimes m} \ket{y} = \sum_{z \in [2^m]} (-1)^{y \cdot z} \ket{z}\, ,
      \intertext{
        and the orthogonality relation,
      }
      \frac{1}{2^m} \sum_{y} (-1)^{y \cdot (z \xor w)} = \delta_{z,w}\, .
    \end{align*}
    Here, $y \cdot z$ denotes the inner product modulo $2$.
  \end{exercise}
